\documentclass[main.tex]{subfiles}
\begin{document}
\ifSubfilesClassLoaded{\begin{center}\small\itshape Part of the GnssProcessor Mathematical Reference \textbullet\ \today\end{center}}{}
\section{Pipeline and architecture diagrams}\label{sec:diagrams}
\newcommand{\Wmat}{\bm{W}}  % observation weight matrix (network design rows)
% ============================================================
%  Pipeline and architecture diagrams (curated TikZ).
%
%  These are hand-authored process diagrams: they show data flow and
%  algorithm phases, which do not fall out of a parser.  Keep them in
%  sync with the code by hand; each box cites the function it maps to.
%  (The machine-generated *reference tables* live in the appendix and
%  are produced by scripts/gen_report_tables.py instead.)
%
%  Shared TikZ styles (proc/io/store/note/flow/feedback) and the
%  \fitwidth helper are defined in main.tex.
% ============================================================

This section shows the control and data flow behind the algebra of the other
chapters: the per-epoch estimation loop (\Cref{fig:epoch-loop}), the shared
geometry kernel behind the three differencing regimes (\Cref{fig:differencing}),
the real-time SSR correction pipeline (\Cref{fig:ssr-pipeline}), the real-time
streaming front-end that feeds it (\Cref{fig:rt-stream}), and the joint
estimator of the PPP-RTK network producer (\Cref{fig:net-flow}).  Each box names
the function or class that implements the step, so the diagrams also serve as a
map into \texttt{src/}.  The diagrams are maintained by hand and simplify the
code: they show the main sequence, not every gate or diagnostic.

% ------------------------------------------------------------
%  Figure 1: per-epoch post-processing loop
% ------------------------------------------------------------
\begin{figure}[tbp]
\centering
\fitwidth{%
\begin{tikzpicture}[node distance=5mm and 12mm]
  \node[io]   (collect)                      {Collect epoch observations\\{\footnotesize decimation $\cdot$ max-duration $\cdot$ gap detection}};
  \node[proc, below=of collect] (geom)       {Geometry models\\{\footnotesize \texttt{PppModel::computeModels}}};
  \node[proc, below=of geom]    (wls)        {WLS clock seeding $\cdot$ gap recovery\\{\footnotesize \texttt{handleWlsSeeding} $\cdot$ \texttt{handleGap}}};
  \node[proc, below=of wls]     (predict)    {Kalman predict\\{\footnotesize \texttt{predict} / \texttt{predictAided}}};
  \node[proc, below=of predict] (slip)       {Cycle-slip detect \& repair\\{\footnotesize \texttt{CycleSlipManager}}};
  \node[proc, below=of slip]    (update)     {Iterative KF update loop\\{\footnotesize outlier gating $\cdot$ pivot handover $\cdot$ measurement rebuild}};
  \node[proc, below=of update]  (relin)      {Geometry relinearisation (IEKF)\\{\footnotesize optional $\cdot$ \texttt{runGeometryRelinearization}}};
  \node[proc, below=of relin]   (gim)        {GIM iono weighting~\textsuperscript{(P)}\\{\footnotesize \texttt{applyGimIonoUpdate}}};
  \node[proc, below=of gim]     (ar)         {Ambiguity resolution\\{\footnotesize MLAMBDA $\cdot$ \texttt{AmbiguityResolver::resolve}}};
  \node[io,   below=of ar]      (log)        {Log epoch + write CSV traces};
  \node[proc, below=8mm of log] (smoother)   {RTS smoother: backward pass~\textsuperscript{(P)}\\{\footnotesize \texttt{Smoother::run}}};

  \node[note, left=10mm of collect, text width=3.4cm] (modes)
        {Before the loop the driver branches on mode:\\
         \texttt{RealTime}~$\to$~\texttt{RealTimeEngine};\\
         baseline~$\to$~DD pipeline;\\
         \texttt{PPP\_RTK\_Network}~$\to$~producer\\
         (or the consumer-replay driver when\\
         \texttt{closed\_loop.}\\\texttt{replay\_products\_path} is set).};

  \foreach \a/\b in {collect/geom, geom/wls, wls/predict, predict/slip,
                     slip/update, update/relin, relin/gim, gim/ar, ar/log}
      \draw[flow] (\a) -- (\b);
  \draw[flow] (log) -- (smoother) node[midway, right, note] {after final epoch};
  \draw[feedback] (log.east) -- ++(1.6,0) |- (collect.east)
      node[pos=0.25, right, note] {next epoch};
\end{tikzpicture}}
\caption{Per-epoch post-processing loop in \texttt{Application::run}.  Solid
arrows give the sequence within an epoch; the dashed arrow is the feedback from
one epoch to the next.  Steps tagged \textsuperscript{(P)} run in
post-processing only; the real-time path (\texttt{RealTimeEngine}) has no GIM
source and no smoother.  The smoother runs once, after the forward pass, and
only when \texttt{run\_smoother} is set.}
\label{fig:epoch-loop}
\end{figure}

% ------------------------------------------------------------
%  Figure 2: differencing strategy hierarchy
% ------------------------------------------------------------
\begin{figure}[tbp]
\centering
\fitwidth{%
\begin{tikzpicture}[node distance=7mm and 6mm]
  \node[io] (raw) {Raw observables\\{\footnotesize rover\,[\,$+$ reference\,]}};
  \node[proc, below=of raw] (kernel)
        {Shared geometry kernel\\{\footnotesize \texttt{PppModel::computeModels}}\\{\footnotesize (observed \& modelled values)}};
  \node[note, right=8mm of kernel, text width=4.9cm]
        {branch chosen by\\\texttt{differencingModeFor(config)}};

  \node[proc, text width=3.4cm, below=18mm of kernel, xshift=-5.4cm] (none)
        {\textbf{None}: PPP\\{\footnotesize no differencing\\(raw uncombined)}};
  \node[proc, text width=3.6cm, below=18mm of kernel] (bssd)
        {\textbf{BetweenSatellite}: PPP-RTK\\{\footnotesize \texttt{BssdTransform::apply}\\SD vs.\ pivot (\texttt{PivotManager})}};
  \node[proc, text width=3.8cm, below=18mm of kernel, xshift=5.6cm] (dd)
        {\textbf{DoubleDifference}: baseline\\{\footnotesize \texttt{BrsdTransform::apply}\\$\to$ \texttt{differenceBrsd}\\(\texttt{DdPivotManager})}};

  \node[proc, below=20mm of bssd] (kf)
        {\texttt{EpochMeasurement} $\to$ \texttt{KalmanFilter::update}};

  \draw[flow] (raw) -- (kernel);
  \draw[flow] (kernel.south) -- (none.north);
  \draw[flow] (kernel.south) -- (bssd.north);
  \draw[flow] (kernel.south) -- (dd.north);
  \draw[flow] (none.south) |- (kf.west);
  \draw[flow] (bssd.south) -- (kf.north);
  \draw[flow] (dd.south)   |- (kf.east);
\end{tikzpicture}}
\caption{The three differencing regimes share one geometry kernel.  The
undifferenced observed and modelled values are computed once.  The regime
returned by \texttt{differencingModeFor} (selected once, when
\texttt{MeasurementBuilderFactory} builds the measurement builder) then applies
no transform (PPP), a between-satellite single difference against a pivot
(PPP-RTK), or a between-receiver$\,\to\,$between-satellite double difference
(baseline).  All three feed the same Kalman update.}
\label{fig:differencing}
\end{figure}

% ------------------------------------------------------------
%  Figure 3: real-time SSR correction pipeline
% ------------------------------------------------------------
\begin{figure}[tbp]
\centering
\fitwidth{%
\begin{tikzpicture}[node distance=7mm and 8mm]
  \node[io] (ntrip) {NTRIP / TCP / file\\{\footnotesize RTCM}};
  \node[io, right=14mm of ntrip] (sbf) {SBF stream\\{\footnotesize rooftop receiver}};
  \node[proc, below=of $(ntrip)!0.5!(sbf)$] (feed)
        {\texttt{feedRtcmFrame} / \texttt{feedSbfFrame}\\{\footnotesize parse \& dispatch}};

  \node[proc, text width=3.0cm, below=10mm of feed, xshift=-5.4cm] (obs)
        {Observations\\{\footnotesize $\to$ \texttt{processEpoch}}};
  \node[proc, text width=3.0cm, below=10mm of feed] (ssrdec)
        {SSR decode\\{\footnotesize legacy 1057--1068\\IGS-SSR 4076}};
  \node[proc, text width=3.2cm, below=10mm of feed, xshift=5.4cm] (ephdec)
        {Broadcast eph.\\{\footnotesize 1019/1020/1042/1044/1046}};

  \node[store, below=10mm of ssrdec] (ssrstore) {\texttt{SsrCorrection}\\\texttt{Store}};
  \node[store, below=10mm of ephdec] (ephstore) {\texttt{LiveEphemeris}\\\texttt{Store}};

  \node[proc, text width=7.2cm, below=24mm of obs] (apply)
        {Apply to undifferenced model\\{\footnotesize orbit (RAC) $\cdot$ clock polynomial $\cdot$ code+phase bias}\\{\footnotesize \texttt{ModelGeometry} $\cdot$ \texttt{GnssCorrections} $\cdot$ \texttt{MeasurementModel}}};
  \node[proc, below=of apply] (reset)
        {Phase-bias discontinuity $\to$ \texttt{resetAmbiguity}};
  \node[proc, below=of reset] (kf)
        {\texttt{KalmanFilter::update} (rebuild loop)};
  \node[proc, below=of kf] (arr)
        {Ambiguity resolution};

  \draw[flow] (ntrip) |- (feed.west);
  \draw[flow] (sbf)   |- (feed.east);
  \draw[flow] (feed.south) -- (obs.north);
  \draw[flow] (feed.south) -- (ssrdec.north);
  \draw[flow] (feed.south) -- (ephdec.north);
  \draw[flow] (ssrdec) -- (ssrstore);
  \draw[flow] (ephdec) -- (ephstore);
  \draw[flow] (obs.south) |- (apply.west);
  \draw[flow] (ssrstore.south) |- (apply.east);
  \draw[flow] (ephstore.south) |- (apply.east);
  \draw[flow] (apply) -- (reset);
  \draw[flow] (reset) -- (kf);
  \draw[flow] (kf) -- (arr);
\end{tikzpicture}}
\caption{Real-time SSR correction pipeline.  The decoder splits the byte stream
three ways: observations go to \texttt{processEpoch}, SSR corrections to
\texttt{SsrCorrectionStore}, and broadcast ephemerides to
\texttt{LiveEphemerisStore}.  The legacy SSR decoder accepts message types
1057--1068 and 1240--1270; the diagram lists the first range only.  The
undifferenced observation model then applies the precise orbit (RAC frame), the
clock polynomial and the code and phase biases.  Between the slip check and the
Kalman update, \texttt{resetOnPhaseBiasDisc} reacts to a change of the
phase-bias discontinuity counter, by default with an ambiguity reset;
\texttt{reinflateOnCodeBiasDisc} and \texttt{reinflateOnClockDisc} follow for
the code-bias and clock counters.  See \texttt{RealTimeEngine::processEpoch}.}
\label{fig:ssr-pipeline}
\end{figure}

% ------------------------------------------------------------
%  Figure 4: real-time streaming front-end (NTRIP -> decode ->
%  sync buffer -> per-epoch update)
% ------------------------------------------------------------
\begin{figure}[tbp]
\centering
\fitwidth{%
\begin{tikzpicture}[node distance=6mm and 10mm]
  \node[io, text width=3.6cm] (ntrip)
        {NTRIP / TCP source\\{\footnotesize \texttt{NtripClient} $\cdot$ \texttt{TcpClient}}};
  \node[proc, text width=4.2cm, below=of ntrip] (reader)
        {Frame extractor\\{\footnotesize \texttt{RtcmStreamReader::feed}\\$\to$~\texttt{nextFrame} (CRC-checked)}};
  \node[proc, text width=4.2cm, below=of reader] (feed)
        {Decode \& dispatch frame\\{\footnotesize \texttt{RealTimeEngine::feedRtcmFrame}}};

  \node[proc, text width=3.0cm, below=10mm of feed, xshift=-5.6cm] (obs)
        {MSM obs $\to$ epoch buffer\\{\footnotesize \texttt{pending\_obs\_}\\(edge-triggered flush)}};
  \node[proc, text width=3.0cm, below=10mm of feed] (ssr)
        {SSR push\\{\footnotesize \texttt{SsrCorrectionStore}}};
  \node[proc, text width=3.0cm, below=10mm of feed, xshift=5.6cm] (eph)
        {Broadcast nav push\\{\footnotesize \texttt{LiveEphemerisStore}}};

  \node[store, text width=2.6cm, below=12mm of obs] (sync)
        {obs$\leftrightarrow$SSR sync buffer\\{\footnotesize \texttt{held\_obs\_}}};
  \node[proc, text width=6.6cm, below=14mm of sync, xshift=2.8cm] (epoch)
        {Per-epoch PPP update\\{\footnotesize \texttt{RealTimeEngine::processEpoch}}\\{\footnotesize WLS seed $\cdot$ predict $\cdot$ slip $\cdot$ KF $\cdot$ AR}};

  \node[note, right=8mm of ssr, text width=3.6cm] (gate)
        {gate (\texttt{enable\_obs\_ssr\_sync}): obs newer than newest SSR content epoch $+$ \texttt{ssr\_max\_age} is held, else processed};

  \draw[flow] (ntrip) -- (reader);
  \draw[flow] (reader) -- (feed);
  \draw[flow] (feed.south) -- (obs.north);
  \draw[flow] (feed.south) -- (ssr.north);
  \draw[flow] (feed.south) -- (eph.north);
  \draw[flow] (obs.south) -- (sync.north);
  \draw[flow] (sync.south) -| (epoch.north);
  \draw[flow] (ssr.south)  |- (epoch.east)
        node[pos=0.25, right, note] {newestSsrEpoch};
  \draw[flow] (eph.south)  |- (epoch.east);
  \draw[feedback] (ssr.west) -- ++(-0.9,0) |- (sync.east)
        node[pos=0.25, above, note] {drainHeldObs};
\end{tikzpicture}}
\caption{Real-time streaming front-end ahead of the SSR correction stage of
\Cref{fig:ssr-pipeline}.  \texttt{RtcmStreamReader} reassembles the raw bytes
from \texttt{NtripClient} or \texttt{TcpClient} into CRC-validated RTCM3 frames,
and \texttt{feedRtcmFrame} decodes and dispatches each frame.  MSM observations
accumulate in \texttt{pending\_obs\_}, which is submitted to
\texttt{processEpoch} when a frame with a new epoch time arrives (the edge
trigger).  SSR corrections and broadcast ephemerides go to their stores.  With
\texttt{enable\_obs\_ssr\_sync} set, the first admission gate of
\texttt{processEpoch} (\texttt{admitBacklogAndSync}) holds an observation epoch
that is newer than the SSR coverage horizon (newest SSR content epoch $+$
\texttt{ssr\_max\_age}).  The held epoch is kept in \texttt{held\_obs\_} (only
the latest one) and resubmitted by \texttt{drainHeldObs} after each SSR push.
An epoch inside the horizon is processed at once.  A hold longer than
\texttt{obs\_ssr\_sync\_max\_hold\_s} (default \SI{30}{s}) latches a degraded
outage state in which epochs pass through without SSR until the stream covers
them again.  The gate applies only with a live SSR orbit and clock source and
never before the first epoch.  See \dcref{sec:rt}{the real-time engine}.}
\label{fig:rt-stream}
\end{figure}

% ------------------------------------------------------------
%  Figure 5: PPP-RTK network producer joint estimator
% ------------------------------------------------------------
\begin{figure}[tbp]
\centering
\fitwidth{%
\begin{tikzpicture}[node distance=7mm and 10mm]
  \node[io, text width=4.6cm] (obs)
        {Multi-station observations\\{\footnotesize \texttt{MultiRinexSync} (by station)}};
  \node[proc, text width=6.2cm, below=of obs] (layout)
        {Register shared sat parameters\\{\footnotesize \texttt{registerClock} $\cdot$ \texttt{registerCodeBias} $\cdot$ \texttt{registerPhaseBias}}};
  \node[proc, text width=6.2cm, below=of layout] (begin)
        {Open epoch arrowhead system\\{\footnotesize \texttt{NetworkEstimator::beginEpoch}}};
  \node[proc, text width=6.2cm, below=of begin] (add)
        {Add per-station design rows\\{\footnotesize \texttt{addStation} ($\jacH_s$, $\jacH_r$, $\Wmat$, $\bm{y}$); one block per station}};
  \node[proc, text width=6.2cm, below=of add] (schur)
        {Schur-reduce local nuisances\\{\footnotesize \texttt{NetworkNormalEquations} arrowhead}};
  \node[proc, text width=6.2cm, below=of schur] (solve)
        {S-basis datum $+$ solve shared\\{\footnotesize \texttt{solve(Datum)} $\cdot$ \texttt{solveSharedConstrained}}};
  \node[proc, text width=6.2cm, below=of solve] (map)
        {Map to solution epoch\\{\footnotesize \texttt{toSolutionEpoch}}};
  \node[io, text width=6.2cm, below=of map] (emit)
        {SSR-style products\\{\footnotesize \texttt{NetworkProductEmitter::emitToStore} (IGS-SSR convention)}};

  \node[note, right=9mm of add, text width=3.6cm]
        {shared cols $=$ (sat, kind, signal); local cols $=$ per-station nuisances (rx clock, ZWD, iono, ambiguities)};
  \node[note, right=9mm of solve, text width=3.6cm]
        {datum pins the gauge freedoms: overall clock, the phase-bias$\leftrightarrow$ambiguity datum, ISB};

  \foreach \a/\b in {obs/layout, layout/begin, begin/add, add/schur,
                     schur/solve, solve/map, map/emit}
      \draw[flow] (\a) -- (\b);
  \draw[feedback] (emit.east) -- ++(1.4,0) |- (obs.east)
        node[pos=0.25, right, note] {next epoch};
\end{tikzpicture}}
\caption{Joint estimator of the PPP-RTK network producer
(\texttt{NetworkEstimator}).  The shared satellite-side columns, one per
(satellite, parameter kind, signal) for the estimable clock, code bias and phase
bias, are registered afresh every epoch (\texttt{reset}, \texttt{registerClock},
\texttt{registerCodeBias}, \texttt{registerPhaseBias}, and
\texttt{registerVtecCoeff} for the SH-VTEC coefficients when
\texttt{network.estimate\_iono\_sh} is on); \texttt{beginEpoch} opens an
arrowhead normal-equation system of that width.  \texttt{addStation} adds one
pre-formed block of local-versus-shared design rows per station, with no
receiver-side differencing.  Schur reduction removes the local nuisances; the
reduced shared system is rank-deficient by the gauge freedoms.  \texttt{solve}
fixes them with an explicit minimum-constraint S-basis datum on a copy of the
reduced system, and \texttt{toSolutionEpoch} maps the estimate to a record that
\texttt{NetworkProductEmitter} writes in the IGS-SSR convention.  The driver
\texttt{Application::runNetwork} runs this core once per epoch over the
\texttt{MultiRinexSync} stations and only orchestrates:
\texttt{networkPassAHead\_} models each station, \texttt{networkAssemblyRegisterColumns\_}
registers the columns and calls \texttt{beginEpoch}, \texttt{networkPassB\_}
calls \texttt{addStation} (all in \texttt{app/ApplicationNetworkAssembly}),
\texttt{networkSolveGlue} applies the datum and solves, \texttt{networkSolPrep}
maps, and \texttt{networkEmitStore} emits.  With
\texttt{network.closed\_loop.enable} set, the emit step also drives a held-out
PPP-RTK consumer on the new products: BSSD by default, undifferenced with
\texttt{closed\_loop.consumer\_undifferenced}.
See \dcref{sec:net}{the network mode} and
\dcref{sec:net-arch}{its producer-architecture section}.}
\label{fig:net-flow}
\end{figure}

\ifSubfilesClassLoaded{\printbibliography}{}
\end{document}
